\documentclass[12pt]{article}
\usepackage[utf8]{inputenc}
\usepackage[T1]{fontenc}
\usepackage{hyperref}
\title{Illyrian Language Resources for K-12 Curriculum Development}
\author{Compiled by Claude Sonnet 5.5 (claude-sonnet-5-5) for Andr\'as Kornai}
\date{}
\begin{document}
\maketitle

\section{Introduction}
Illyrian (ISO 639-3 xil) is an extinct, unclassified language of the ancient Balkans. Nothing was written in it. It is known only from a few glosses recorded by Greek and Roman authors and from personal and place names. Scholars disagree about its extent and relationships, and some argue that ``Illyrian'' is better understood as a historical-geographic label than as a single language. Any ``native language'' curriculum is therefore not possible. A heritage curriculum can only teach about the language and the evidence for it.

\section{Primary Educational Resources}
I found no primers, textbooks, or school programmes for Illyrian.

\section{Linguistic and Reference Materials}
\begin{itemize}
\item Illyrian language (Wikipedia): overview of the evidence, glosses, onomastics, and proposed relationships. \url{https://en.wikipedia.org/wiki/Illyrian_language}
\item \emph{Lucida Intervalla} no.\ 51 (2022): journal issue with a paper on Palaeo-Balkan onomastics. \url{https://www.f.bg.ac.rs/lucidaintervalla/issues/51(2022).pdf}
\item Sa\v{v}i\'c, \emph{Lucida Intervalla} 51, in the Ligorio paper series. \url{https://www.klasicnenauke.rs/wp-content/uploads/2023/05/LI51-Savic-Ligorio.pdf}
\item Article in \emph{Rhesis} on Illyrian linguistic evidence. \url{https://ojs.unica.it/index.php/rhesis/article/download/5594/5463}
\end{itemize}

\section{Religious and Community Resources}
No religious texts or community language organisations were found. Illyrian has no surviving texts and no known speakers.

\section{Digital and Audio Resources}
I found no audio or digital language resources for Illyrian.

\section{Curriculum Development Context}
A heritage unit could cover the ancient Illyrians, how a language can be reconstructed from names and glosses, and the debated link with Albanian. The sources above are scholarly and written for adults, so they would need adaptation for K-12 use.
\end{document}